
% PROBLEMIDENTIFIKATION VERSION 0.0 %

I det moderne, krise- og krigsprægede samfund står vi over for en række udfordringer, som påvirker vores liv, livskvalitet og fremtidige muligheder. De mest centrale problemområder identificeres som følgende:

\begin{enumerate}
    \item Økonomi: Stigende priser, lavere købekraft og en generel usikkerhed omkring privatøkonomien.
    \item Miljø og klima: Behovet for bæredygtige løsninger inden for energi, affaldshåndtering og reduktion af $CO_2$-udslip.
    \item Geopolitisk ustabilitet og krig: Internationale spændinger og konflikter truer den globale sikkerhed og påvirker energipriser, inflation og forsyningskæder.
\end{enumerate}

Ovenstående problemområder præger i høj grad samfundet i dag og har direkte konsekvenser for den enkelte borger. Det er derfor vigtigt, at vi som samfund arbejder mod løsninger, der kan forbedre eksempelvis privatøkonomien og dermed øge livskvaliteten.

Inden for det økonomiske område handler problemstillingerne blandt andet om høje priser og, hvordan man kan modvirke disse. Der er fokus på, hvordan den almene dansker kan effektivisere sin privatøkonomi og opnå besparelser med minimale ændringer i hverdagen.

Miljømæssigt omfatter problemstillingerne blandt andet forurening, affaldssortering og reduktion af det samlede miljøaftryk. Klimamæssigt er fokus på at mindske $CO_2$-udslip og forbedre klimaaftrykket, ligeledes med så små ændringer i hverdagen som muligt.

I forhold til krig og konflikter kan problemstillingerne omhandle, hvordan man reducerer tab af menneskeliv, forbedrer nødhjælp til udsatte områder, eksempelvis gennem levering af feltrationer, samt hvordan dette kan ske uden at udsætte yderligere parter for risiko.

På baggrund af de økonomiske problemstillinger opstår idéen om et produkt, der kan hjælpe forbrugere med at få bedre overblik over deres abonnementer. Produktet kan eksempelvis være en enhed eller applikation, der analyserer netværkstrafik fra brugerens router og registrerer, hvilke hjemmesider og apps der anvendes. På den baggrund kan der genereres statistik over brugen af forskellige abonnementstjenester, hvilket gør det muligt at identificere unødvendige abonnementer og dermed reducere udgifter.


% AFGRÆNSNING VERSION 0.0 %
\subsection{Afgrænsning}
Vi har valgt at Afgrænse vores projekt således at vi kun fokuserer på digitale abonnementer, som ikke indeholder personfølsomme data, såsom Bank data og abonnementer der er tilkoblet meget personfølsomme informationer, som kan ødelægge vigtig infrastruktur for individet, således der ikke er nogle funktioner af vores produkt der overskrider persondatalovgivningen

% PROBLEMFORMULERING VERSION 0.0 %
\subsection{Problemformulering}
Hvordan kan vi ved brug af moderne mobilteknologi bekæmpe fænomenet Subscription Fatigue hos danske forbrugere, således de opnår et bedre økonomisk overblik og reducerer mentalt stress?

For at besvare dette vil vi undersøge følgende:
\begin{itemize}
    \item Hvilke økonomiske og psykologiske effekter ligger bag Subscription Fatigue og hvordan påvirker det vores målgruppes privatøkonomi?
    \item Hvordan kan vi designe et produkt og udvikle den således den kan give brugeren en automatiseret og overskuelig analyse af deres faktiske forbrug af abonnementstjenester?
    \item Hvordan kan produktet markedsføres effektivt gennem Growth Hacking-principper, og hvilke miljømæssige overvejelser bør man gøre sig i forbindelse med driften af de bagvedliggende datacentre?
\end{itemize}




\subsection{Projektstyringsværktøj (meta)}
Hele projektet kan findes via platformen GitHub, hvor vi har oprettet et repository til projekt. På GitHub kan det ses, hvordan projektet har udviklet sig over tid, samt rapportens rene \LaTeX-kode og produktets kildekode.\\ 

Linket til repository: \url{https://github.com/ZBC-Slagelse-HTX-X/teknologi_a_projekt}\\
For vejledning åben linktet og følg instruksionerne.

Det endelige produkt er afleveret som APK, på et USB-stik, men kan også hentes via Githubs release-funktionalitet. Den afleverede version er alpha-v.1.0.0. 

Det er vigtigt at notere sig at APK-filer kun kan køres på android-enheder, herunder android-emulatorere. Appen er testet på 3 mobiltelefoner (S25 Ultra, S22 Ultra samt google pixel 10 pro), det anbefales derfor at man IKKE benytter appen på en fold-telefon, da det ikke har været muligt at teste.